% \chapter{Diagonalisierung von $\dul{h}$ bei beliebiger Füllung}~\label{sec:anhang_a}

% Die Matrix $\dul{h}$~\eqref{eq:h_nothalf} besitzt die Blockstruktur
% \begin{align}
%     \dul{h} &= \begin{pmatrix}
%         \dul{M} & \Delta \onematrix \\
%         \Delta \onematrix & -\dul{M}
%     \end{pmatrix} & \text{mit} & & \dul{M} &= \begin{pmatrix}
%         \mu & -f^*(\vk) \\
%         -f(\vk) & \mu
%     \end{pmatrix} \, .
% \end{align}
% Diese spezielle Struktur kann nun ausgenutzt werden, um das Eigenwertproblem einer $4 \times 4$-Matrix auf das einiger $2 \times 2$-Matrizen zu reduzieren.
% Im Folgenden werden Matrizen der Übersichtlichkeit halber fett markiert, also $\dul{A} \rightarrow \mb{A}$.

% Es sei $\mUM$ die unitäre Trafo, die $\mM$ diagonalisiert, für die also
% \begin{align}
%     \mM_\diag &= \diag(\lambda_+, \lambda_-) = \mUM^\dagg  \mM \, \mUM
% \end{align}
% gilt.
% Dann kann $\mh$ mit $\mW = \diag(\mUM, \mUM)$ wiefolgt transformiert werden:
% \begin{align}
%     \mhp &= \mW^\dagg \mh \, \mW \\
%     &= \begin{pmatrix}
%         \mUM^\dagg & 0 \\ 0 & \mUM^\dagg
%     \end{pmatrix} \cdot \begin{pmatrix}
%         \mM & \Delta \mone \\ \Delta \mone & -\mM
%     \end{pmatrix} \cdot \begin{pmatrix}
%         \mUM & 0 \\ 0 & \mUM
%     \end{pmatrix} \\
%     &= \begin{pmatrix}
%         \mUM^\dagg \mM \, \mUM & \Delta \mone \\ \Delta \mone &-\mUM^\dagg \mM \, \mUM 
%     \end{pmatrix} \\
%     &= \begin{pmatrix}
%         \lambda_+ & 0 & \Delta & 0 \\
%         0 & \lambda_- & 0 & \Delta \\
%         \Delta & 0 & -\lambda_+ & 0 \\
%         0 & \Delta & 0 & -\lambda_-
%     \end{pmatrix} \, .
% \end{align} 
% Diese transformierte Matrix $\mhp$ kann wiederum durch gleichzeitigem Tausch der zweiten und dritten Zeile und Spalte in eine Blockstruktur gebracht werden. 
% Dieser Tausch wird formell als weitere unitäre Trafo 
% \begin{align}
%     \mhpp &= \mP^\dagg \mhp \mP = \begin{pmatrix}
%         \lambda_+ & \Delta & 0 & 0 \\
%         \Delta & - \lambda_+ & 0 & 0 \\
%         0 & 0 & \lambda_- & \Delta \\
%         0 & 0 & \Delta & -\lambda_-
%     \end{pmatrix} = \begin{pmatrix}
%         \mBp & 0 \\
%         0 & \mBm
%     \end{pmatrix}
% \end{align}
% mit der Permutationsmatrix 
% \begin{align}
%     \mP &= \begin{pmatrix}
%         1 & 0 & 0 & 0 \\
%         0 & 0 & 1 & 0 \\
%         0 & 1 & 0 & 0 \\
%         0 & 0 & 0 & 1
%     \end{pmatrix}
% \end{align}
% beschrieben.

% % In dieser Form können nun $\mBpm$ unabhängig voneinander mit $\mUpm$ diagonalisiert werden und damit $\mhpp$ ergibt sich schließlich 

% Diagonalisieren die unitären Transformationen $\mUpm$ nun die beiden Blockmatrizen $\mBpm$, wird $\mhpp$ mit $\mWp = \diag(\mBp, \mBm)$ diagonal. 
% Es gilt 
% \begin{align}
%     \mhpp_\diag &= \mWp^\dagg \mhpp \mWp = \diag(\mUp^\dagg \mBp \mUp, \mUm^\dagg \mBm \mUm) = \diag(E_+^{(1)}, E_+^{(2)}, E_-^{(1)}, E_-^{(2)}) \, ,
% \end{align}
% mit den Eigenwerten $E_{\pm}^{(1/2)}$ von $\mBpm$.

% Zusammengefasst gilt also 
% \begin{align}
%     \mh_\diag = \mWp^\dagg \mP^\dagg \mW^\dagg \mh \mW \mP \mWp = (\mW \mP \mWp)^\dagg \mh \, (\mW \mP \mWp) = \diag(E_+^{(1)}, E_+^{(2)}, E_-^{(1)}, E_-^{(2)}) \, ,
% \end{align}
% also ist die gesuchte Trafo $\mU = \mW \mP \mWp$.

% Es ist leicht nachgerechnet, dass 
% \begin{align}
%     \mUM &= \frac{1}{\sqrt{2}}\begin{pmatrix}
%         -e^{-\im \varphi} & e^{-\im \varphi} \\ 1 & 1
%     \end{pmatrix} & \text{mit} & & \lambda_\pm &= \mu \pm \abs{f(\vk)} 
% \end{align}
% die Matrix $\mM$ diagonalisiert, während $\mBpm$ durch die Matrizen 
% \begin{align}
%     \mUp &= \frac{1}{\sqrt{2}}\begin{pmatrix}
%         \frac{\Delta}{N_1} & \frac{\Delta}{N_2} \\
%         \frac{E_+ - \lambda_+}{N_1} & \frac{-E_+ - \lambda_+}{N_2}
%     \end{pmatrix} & \text{mit} & & N_{1,2} &= \sqrt{E_+(E_+ \mp \lambda_+)}
% \end{align}
% und 
% \begin{align}
%     \mUm &= \frac{1}{\sqrt{2}}\begin{pmatrix}
%         \frac{\Delta}{N_3} & \frac{\Delta}{N_4} \\
%         \frac{E_- - \lambda_-}{N_3} & \frac{-E_- - \lambda_-}{N_4}
%     \end{pmatrix} & \text{mit} & & N_{3,4} &= \sqrt{E_-(E_- \mp \lambda_-)}
% \end{align}
% und $E_\pm = \sqrt{\lambda_\pm^2 + \Delta^2}$ in ihre Eigenwerte $E_{+}^{(1/2)} = \pm E_+$, sowie $E_-^{(1/2)} = \pm E_-$ zerlegt werden.




\chapter{Diagonalisierung von $\dul{h}$ bei beliebiger Füllung}\label{sec:anhang_a}

Die Matrix $\dul{h}$~\eqref{eq:h_nothalf} besitzt die Blockstruktur
\begin{align}
    \dul{h} &= \begin{pmatrix}
        \dul{M} & \Delta \onematrix \\
        \Delta \onematrix & -\dul{M}
    \end{pmatrix} & \text{mit} & & \dul{M} &= \begin{pmatrix}
        \mu & -f^*(\vk) \\
        -f(\vk) & \mu
    \end{pmatrix} 
\end{align}
und es wird die Matrix $\dul{U}$ gesucht, die $\dul{h}$ diagonalisiert.
Diese Struktur lässt sich gezielt ausnutzen, um das Eigenwertproblem der $4\times4$-Matrix $\dul h$ auf mehrere $2\times2$-Probleme zurückzuführen, die sich unabhängig voneinander lösen lassen. Der Übersichtlichkeit halber werden Matrizen im Folgenden fett markiert, also $\dul{A} \rightarrow \mb{A}$.

% \paragraph{Schritt 1 -- Diagonalisierung von $\mM$.}
Es sei $\mUM$ die unitäre Transformation, die $\mM$ diagonalisiert, für die also
\begin{align}
    \mM_\diag &= \diag(\lambda_+, \lambda_-) = \mUM^\dagg  \mM \, \mUM
\end{align}
gilt. 
Wendet man diese Transformation mit $\mW = \diag(\mUM, \mUM)$ auf die Diagonalblöcke $\mM$ an, so bleibt die Blockstruktur von $\mh$ erhalten, während die Diagonalblöcke selbst diagonal werden:
\begin{align}
    \mhp &= \mW^\dagg \mh \, \mW \\
    &= \begin{pmatrix}
        \mUM^\dagg & 0 \\ 0 & \mUM^\dagg
    \end{pmatrix} \cdot \begin{pmatrix}
        \mM & \Delta \mone \\ \Delta \mone & -\mM
    \end{pmatrix} \cdot \begin{pmatrix}
        \mUM & 0 \\ 0 & \mUM
    \end{pmatrix} \\
    &= \begin{pmatrix}
        \mUM^\dagg \mM \, \mUM & \Delta \mone \\ \Delta \mone &-\mUM^\dagg \mM \, \mUM 
    \end{pmatrix} \\
    &= \begin{pmatrix}
        \lambda_+ & 0 & \Delta & 0 \\
        0 & \lambda_- & 0 & \Delta \\
        \Delta & 0 & -\lambda_+ & 0 \\
        0 & \Delta & 0 & -\lambda_-
    \end{pmatrix} \, .
\end{align}

% \paragraph{Schritt 2 -- Umsortierung in echte Blockdiagonalform.}
% In $\mhp$ sind nun zwar die Diagonalblöcke diagonal, über die $\Delta$-Terme koppeln aber weiterhin Zustand $1$ mit $3$ sowie Zustand $2$ mit $4$. 
In $\mhp$ sind nun zwar die Diagonalblöcke diagonal, aber aufgrund der $\Delta$-Terme ergibt sich noch keine echte diagonale Form.
Vertauscht man allerdings gleichzeitig die zweite und dritte Zeile bzw.\ Spalte, zerfällt $\mhp$ in eine blockdiagonale Form. 
Dieser Tausch entspricht formal einer weiteren unitären Transformation mit der Permutationsmatrix
\begin{align}
    \mP &= \begin{pmatrix}
        1 & 0 & 0 & 0 \\
        0 & 0 & 1 & 0 \\
        0 & 1 & 0 & 0 \\
        0 & 0 & 0 & 1
    \end{pmatrix} \, ,
\end{align}
und liefert
\begin{align}
    \mhpp &= \mP^\dagg \mhp \mP = \begin{pmatrix}
        \lambda_+ & \Delta & 0 & 0 \\
        \Delta & - \lambda_+ & 0 & 0 \\
        0 & 0 & \lambda_- & \Delta \\
        0 & 0 & \Delta & -\lambda_-
    \end{pmatrix} = \begin{pmatrix}
        \mBp & 0 \\
        0 & \mBm
    \end{pmatrix} \, .
\end{align}

% \paragraph{Schritt 3 -- Diagonalisierung der Blöcke.}
Die Blockmatrizen $\mBp$ und $\mBm$ lassen sich nun unabhängig voneinander diagonalisieren. Werden $\mBpm$ durch die unitären Transformationen $\mUpm$ diagonalisiert, so wird $\mhpp$ mit $\mWp = \diag(\mUp, \mUm)$ diagonal:
\begin{align}
    \mhpp_\diag &= \mWp^\dagg \mhpp \mWp = \diag(\mUp^\dagg \mBp \mUp, \mUm^\dagg \mBm \mUm) = \diag(E_+^{(1)}, E_+^{(2)}, E_-^{(1)}, E_-^{(2)}) \, ,
\end{align}
mit den Eigenwerten $E_{\pm}^{(1/2)}$ von $\mBpm$.

% \paragraph{Ergebnis.}
Da jede der drei Transformationen $\mW$, $\mP$ und $\mWp$ unitär ist, ist auch ihr Produkt unitär, und es gilt insgesamt
\begin{align}
    \mh_\diag = \mWp^\dagg \mP^\dagg \mW^\dagg \mh \mW \mP \mWp = (\mW \mP \mWp)^\dagg \mh \, (\mW \mP \mWp) = \diag(E_+^{(1)}, E_+^{(2)}, E_-^{(1)}, E_-^{(2)}) \, ,
\end{align}
die gesuchte Transformation ist also $\mU = \mW \mP \mWp$.

Durch direktes Nachrechnen findet man für die einzelnen Bausteine explizit
\begin{align}
    \mUM &= \frac{1}{\sqrt{2}}\begin{pmatrix}
        -e^{-\im \varphi} & e^{-\im \varphi} \\ 1 & 1
    \end{pmatrix} & \text{mit} & & e^{-\im \varphi} = \frac{f^*(\vk)}{\abs{f(\vk)}} \, ,
\end{align}
die $\mM$ in die Eigenwerte $\lambda_\pm = \mu \pm \abs{f(\vk)}$ zerlegt, 
sowie die Matrizen $\mUpm$, die $\mBpm$ in ihre Eigenwerte $E_{+}^{(1/2)} = \pm E_+$ bzw.\ $E_-^{(1/2)} = \pm E_-$ (mit $E_\pm = \sqrt{\lambda_\pm^2 + \Delta^2}$) zerlegen:
\begin{align}
    \mUp &= \frac{1}{\sqrt{2}}\begin{pmatrix}
        \frac{\Delta}{N_1} & \frac{\Delta}{N_2} \\
        \frac{E_+ - \lambda_+}{N_1} & \frac{-E_+ - \lambda_+}{N_2}
    \end{pmatrix} & \text{mit} & & N_{1,2} &= \sqrt{E_+(E_+ \mp \lambda_+)} \, ,
\end{align}
\begin{align}
    \mUm &= \frac{1}{\sqrt{2}}\begin{pmatrix}
        \frac{\Delta}{N_3} & \frac{\Delta}{N_4} \\
        \frac{E_- - \lambda_-}{N_3} & \frac{-E_- - \lambda_-}{N_4}
    \end{pmatrix} & \text{mit} & & N_{3,4} &= \sqrt{E_-(E_- \mp \lambda_-)} \, .
\end{align}

Werden diese Matrizen dann gemäß $\mU = \mW \mP \mWp$ ausmultipliziert ergibt sich schließlich, die in Gleichung~\eqref{eq:U_nothalf} angegebene unitäre Gesamttrafo 
\begin{align}
    \mU = \frac{1}{2} \begin{pmatrix}
        -\frac{\Delta}{N_1} e^{-\im \varphi} \quad & -\frac{\Delta}{N_2} e^{-\im \varphi} \quad & -\frac{\Delta}{N_3} e^{-\im \varphi} \quad & -\frac{\Delta}{N_4} e^{-\im \varphi} \\[0.3em]
        \frac{\Delta}{N_1} & \frac{\Delta}{N_2} & \frac{\Delta}{N_3} & \frac{\Delta}{N_4} \\[0.3em]
        \left(\frac{-E_+ + \lambda_+}{N_1}\right)e^{-\im \varphi} \quad & \left(\frac{E_+ + \lambda_+}{N_2}\right)e^{-\im \varphi} \quad & \left(\frac{E_- - \lambda_-}{N_3}\right)e^{-\im \varphi} \quad & \left(\frac{-E_- - \lambda_-}{N_4}\right)e^{-\im \varphi} \\[0.3em]
        \frac{E_+ - \lambda_+}{N_1} & \frac{-E_+ - \lambda_+}{N_2} & \frac{E_- - \lambda_-}{N_3} & \frac{-E_- - \lambda_-}{N_4}
    \end{pmatrix}
\end{align}
mit der Diagonalmatrix
\begin{align}
    \mh_\diag = \text{diag}(E_+, -E_+, E_-, -E_-)
\end{align}
aus Gleichung~\eqref{eq:h_diag_nothalf}.